\begin{agradecimentos}

Out of respect for those I wish to thank, these acknowledgements are written in Portuguese, my native language.

Não seria possível começar esses agradecimentos sem mencionar meus pais, Patricia e Jeferson, e meu irmão Gabriel, que fizeram o seu melhor para me apoiar ao longo dos anos.
Sei que não foi fácil acompanhar meus passos de longe, mas sou muito grato pela confiança que depositaram em mim.
Agradeço também ao Leandro e à Kênnia por acompanhá-los durante o dia a dia; vocês foram essenciais.
Tudo por vocês ainda é pouco.

Ao Caio e ao Tulio, meus queridos amigos e companheiros de apartamento: carregarei com muito carinho todas as memórias e registros que fizemos juntos.
Obrigado por todas as gargalhadas e por tudo que me ensinaram.
E, claro, peço desculpas por qualquer coisa!

Aos meus amigos Ezequiel, Carol, Pablo, Isaque, Juliano, Guilherme e Jairon, obrigado por todas as interações especiais que tivemos juntos.

Aos professores Yuri Saporito, Rodrigo Targino, Diego Mesquita, Hugo de la Cruz e Marcelo Sant'Anna, deixo meu agradecimento por todos os conselhos, pelo conhecimento compartilhado e, sobretudo, pelo suporte nas decisões mais importantes da minha trajetória.

Aos meus co-autores em outro trabalho, Guilherme T. Goedert e Daniel Csillag, obrigado pelo suporte e pela colaboração que tornaram possível construir um trabalho de qualidade num espaço com tanta visibilidade como o NeurIPS.

À Cássia Pessanha, agradeço o carinho e o cuidado que teve comigo, sobretudo na ausência de familiares por perto.
Aos professores César Camacho e Maria Izabel Camacho, agradeço a oportunidade de ingressar na EMAp através do FGV CDMC, e ao Professor César, também, por toda a orientação e apoio à frente da direção da escola.
Aos funcionários Ivone, Mônica, Marcio, Jorge, Ingrid, Roseni, Thais, João, Isabel e Cristiane, obrigado pela dedicação ao funcionamento da nossa instituição.
Aos meus alunos de monitoria e aos professores responsáveis pelas disciplinas, obrigado pela confiança e pelo aprendizado mútuo.
Ensinar é, sem dúvida, uma das melhores formas de aprender.
Vocês são minha família, e podem contar comigo sempre.
Serei eternamente grato!

Por fim, ao meu co-orientador Flávio Gonçalves, que me instigou a reflexões fundamentais na minha jornada, e ao meu orientador Luiz Max Carvalho, que nunca permitiu que eu questionasse meu potencial e esteve ao meu lado em todos os momentos: levarei para a vida todas as boas (e as não tão boas) práticas que aprendi com vocês, que servem de motivação para o professor e pesquisador que almejo me tornar.
Obrigado pela confiança e pela oportunidade de trabalhar com vocês.
O trabalho não acaba aqui.

\end{agradecimentos}